\section{Empirical Design}

\subsection{Regression Framework}

To quantify the predictive content of microstructure-based measures, we estimate cross-sectional regressions of post-earnings returns on both the constructed pressure score and the extracted principal components.

The baseline specification is:

\begin{equation}
R_{i,t \rightarrow t+H} = \alpha + \beta_1 \cdot \text{PressureScore}_{i,t} + \gamma' \mathbf{X}_{i,t} + \epsilon_{i,t}
\end{equation}

where $R_{i,t \rightarrow t+H}$ denotes cumulative returns over horizon $H \in \{5,20\}$ trading days, and $\mathbf{X}_{i,t}$ includes control variables such as size, momentum, and volatility.

To further investigate the latent structure, we estimate:

\begin{equation}
R_{i,t \rightarrow t+H} = \alpha + \sum_{k=1}^{K} \beta_k \cdot PC_{k,i,t} + \epsilon_{i,t}
\end{equation}

where $PC_k$ are the principal components derived from the microstructure feature space.

\paragraph{Interpretation}

This framework allows us to test whether:
\begin{itemize}
    \item The economically constructed pressure score has predictive power for returns
    \item The latent factors extracted via PCA carry incremental explanatory power
    \item The mapping between statistical factors and economic mechanisms (information vs liquidity) is empirically supported
\end{itemize}

Standard errors are computed using heteroskedasticity-robust estimators.

Additionally, we further estimate interaction specifications to test whether the effect of pressure depends on the magnitude of earnings news:

\begin{equation}
R_{i,t \rightarrow t+H} = \alpha + \beta_1 \cdot \text{PressureScore}_{i,t} + \beta_2 \cdot \text{SUE}_{i,t} + \beta_3 \cdot (\text{PressureScore}_{i,t} \times \text{SUE}_{i,t}) + \epsilon_{i,t}
\end{equation}
